\texttt{journalctl} logt informatie op basis van prioriteit en facility. De prioriteit geeft aan hoe ernstig een melding is en de facility geeft aan waar het bericht bij hoort.

Als we meldingen met een bepaalde prioriteit willen zien dan moeten we de \texttt{-p} optie gebruiken. Hier kunnen we de naam van een prioriteit of het nummer van de prioriteit gebruiken. Prioriteiten worden genummerd van 0 tot 7 met 0 als hoogste prioriteit. De prioriteiten met hun naam staat in de volgende lijst met Severity beschrijvingen.

\begin{flushleft}
\begin{table}[h!]
\centering
\begin{tabularx}{\textwidth}{ |c|c|c|X| }
\hline
Code  & Severity & Keyword & Beschrijving \\
\hline
0 & Emergency & emerg & System is onbruikbaar, ernstige Kernel BUG, systemd dumped core. Alleen voor de kernel, niet voor applicaties. \\
\hline
1 & Alert & alert & Moet meteen opgelost worden, vaak een vitaal subsysteem dat niet meer werkt, kans op data verlies. \\
\hline
2 & Critical & crit & Kritieke conditie, crashes, coredumps, etc. \\
\hline
3 & Error & err & Error condities, niet fatale errors \\
\hline
4 & Warning & warning & Geeft aan dat er een error zou kunnen optreden als er geen actie wordt ondernomen. Komt bijvoorbeeld voor bij het vollopen van een disk \\
\hline
5 & Notice & notice & Gebeurtenissen die ongewoon zijn, maar die geen error opleveren, misschien wel iets om naar te kijken \\
\hline
6 & Informational & info & Berichten die geen actie behoeven, maar wel gelogd willen worden \\
\hline
7 & Debug & debug & Berichten waarbij vaak eerst de debug optie aangezet moet worden \\
\hline
\end{tabularx}
\caption{Severity beschrijving}
\label{table:facility}
\end{table}
\end{flushleft}

We kunnen ook een range van prioriteiten opgeven
\begin{lstlisting}[language=bash]
$ sudo journalctl -p 0..3
\end{lstlisting}
Alle prioriteiten die genoemd zijn worden weergegeven, dit is dus in dit voorbeeld 0 tot en met 3.

Voor de faciliteit waartoe een bericht hoort wordt er gebruik gemaakt van de facility zoals deze beschreven is in RFC 5424 6.2.1. De verschillende faciliteiten zijn beschreven in de tabel met Facility beschrijvingen.

\begin{flushleft}
\begin{table}[h!]
\centering
\begin{tabularx}{\textwidth}{ |c|c|X|X| }
\hline
	Code &
	Keyword &
	Beschrijving &
	Informatie \\
\hline
0 & kern & Kernel messages & \\
\hline
1 & user & User-level messages & \\
\hline
2 & mail & Mail system & Oude POSIX faciliteit, soms nog gebruikt, gebruik -u met de unit-name van de mailserver \\
\hline
3 & daemon & System daemons & Alle daemons (services), inclusief systemd en kernel daemons \\
\hline
4 & auth & Security/authorization messages & Zie ook facility 10 \\
\hline
5 & syslog & Berichten van syslogd als deze wordt gebruikt & Wordt niet gebruikt voor het loggen in systemd, zie hiervoor facility 3 \\
\hline
6 & lpr & Line printer subsysteem & oud subsysteem, gebruik -u met de unit-name van de printerserver (cups) \\
\hline
7 & news & Network news subsysteem & oud subsysteem \\
\hline
8 & uucp & UUCP subsysteem & oud subsysteem \\
\hline
9 &      & Clock daemon & systemd-timesyncd \\
\hline
10 & authpriv & Security/authorization messages & Zie ook facility 4 \\
\hline
11 & ftp & FTP daemon & Gebruik -u met de unit-naam van de FTP-server \\
\hline
12 &     & NTP subsystem & Zie facility 9 of gebruik -u met de unit-naam van de NTP-server \\
\hline
13 &     & Log audit & \\
\hline
14 &     & Log alert & \\
\hline
15 & cron & Scheduling daemon & \\
\hline
16 & local0 & Local use 0 (local0) & \\
\hline
17 & local1 & Local use 1 (local1) & \\
\hline
18 & local2 & Local use 2 (local2) & \\
\hline
19 & local3 & Local use 3 (local3) & \\
\hline
20 & local4 & Local use 4 (local4) & \\
\hline
21 & local5 & Local use 5 (local5) & \\
\hline
22 & local6 & Local use 6 (local6) & \\
\hline
23 & local7 & Local use 7 (local7) & \\
\hline
\end{tabularx}
\caption{Facility beschrijving}
\label{table:facility}
\end{table}
\end{flushleft}


Nuttige facilities om te volgen zijn 0, 1, 3, 4, 9, 10, 15.

